\section{8. Diseño de la interfaz de
usuario}\label{diseuxf1o-de-la-interfaz-de-usuario}

\textcolor{red}{\textbf{Pendiente (corte Review 5, Sprints 4 y 5).} La API del flujo núcleo ya existe
(sección 9); el prototipo de alta fidelidad, la evaluación heurística y la prueba con usuarios se hacen
sobre los wireframes de 4.10 una vez validados con el cliente.}

\subsection{8.1 Prototipo de alta
fidelidad}\label{prototipo-de-alta-fidelidad}

\begin{quote}
\emph{\textbf{Instrucción:} Enlace al prototipo navegable y capturas del
flujo principal y de dos flujos adicionales. Expliquen las decisiones de
diseño según el contexto de uso de cada tipo de usuario (dispositivo,
conectividad, experiencia).}
\end{quote}

{[}Enlace y figuras.{]}

\subsection{8.2 Accesibilidad}\label{accesibilidad}

\begin{quote}
\emph{\textbf{Instrucción:} Qué criterios de WCAG 2.2 se consideraron
(contraste, tamaño de objetivos táctiles, textos alternativos,
navegación por teclado) y cómo se verificaron.}
\end{quote}

{[}Escriba aquí.{]}

\subsection{8.3 Evaluación heurística
recibida}\label{evaluaciuxf3n-heuruxedstica-recibida}

\begin{quote}
\emph{\textbf{Instrucción:} Hallazgos que reportó el otro equipo con las
10 heurísticas de Nielsen, y cómo se atendieron.}
\end{quote}

{\def\LTcaptype{none} % do not increment counter
\begin{longtable}[]{@{}
  >{\raggedright\arraybackslash}p{(\linewidth - 8\tabcolsep) * \real{0.1174}}
  >{\raggedright\arraybackslash}p{(\linewidth - 8\tabcolsep) * \real{0.2138}}
  >{\raggedright\arraybackslash}p{(\linewidth - 8\tabcolsep) * \real{0.3056}}
  >{\raggedright\arraybackslash}p{(\linewidth - 8\tabcolsep) * \real{0.1282}}
  >{\raggedright\arraybackslash}p{(\linewidth - 8\tabcolsep) * \real{0.2350}}@{}}
\toprule\noalign{}
\begin{minipage}[b]{\linewidth}\raggedright
\textbf{Issue}
\end{minipage} & \begin{minipage}[b]{\linewidth}\raggedright
\textbf{Heurística}
\end{minipage} & \begin{minipage}[b]{\linewidth}\raggedright
\textbf{Hallazgo}
\end{minipage} & \begin{minipage}[b]{\linewidth}\raggedright
\textbf{Severidad}
\end{minipage} & \begin{minipage}[b]{\linewidth}\raggedright
\textbf{Acción tomada}
\end{minipage} \\
\midrule\noalign{}
\endhead
\bottomrule\noalign{}
\endlastfoot
& & & & \\
& & & & \\
& & & & \\
\end{longtable}
}

\subsection{8.4 Prueba con usuarios y
SUS}\label{prueba-con-usuarios-y-sus}

\begin{quote}
\emph{\textbf{Instrucción:}} \emph{Prueba con} \emph{\textbf{3 a 5
personas}} \emph{ajenas al equipo:} \emph{tareas que realizaron,
problemas observados, puntaje} \emph{promedio del SUS con su cálculo y
cambios derivados.}
\end{quote}

{\def\LTcaptype{none} % do not increment counter
\begin{longtable}[]{@{}
  >{\raggedright\arraybackslash}p{(\linewidth - 10\tabcolsep) * \real{0.2138}}
  >{\raggedright\arraybackslash}p{(\linewidth - 10\tabcolsep) * \real{0.1816}}
  >{\raggedright\arraybackslash}p{(\linewidth - 10\tabcolsep) * \real{0.1174}}
  >{\raggedright\arraybackslash}p{(\linewidth - 10\tabcolsep) * \real{0.1069}}
  >{\raggedright\arraybackslash}p{(\linewidth - 10\tabcolsep) * \real{0.2521}}
  >{\raggedright\arraybackslash}p{(\linewidth - 10\tabcolsep) * \real{0.1281}}@{}}
\toprule\noalign{}
\begin{minipage}[b]{\linewidth}\raggedright
\textbf{Participante (perfil, sin nombre)}
\end{minipage} & \begin{minipage}[b]{\linewidth}\raggedright
\textbf{Tarea}
\end{minipage} & \begin{minipage}[b]{\linewidth}\raggedright
\textbf{¿La completó?}
\end{minipage} & \begin{minipage}[b]{\linewidth}\raggedright
\textbf{Tiempo}
\end{minipage} & \begin{minipage}[b]{\linewidth}\raggedright
\textbf{Problemas observados}
\end{minipage} & \begin{minipage}[b]{\linewidth}\raggedright
\textbf{Puntaje SUS}
\end{minipage} \\
\midrule\noalign{}
\endhead
\bottomrule\noalign{}
\endlastfoot
& & & & & \\
& & & & & \\
& & & & & \\
\end{longtable}
}

\textbf{Puntaje SUS promedio:} {[} {]} · \textbf{Cambios realizados a
partir de la prueba:} {[} {]}
